\documentclass[10pt,twocolumn]{article}

\usepackage[
  letterpaper,
  top=0.66in,
  bottom=0.63in,
  left=0.66in,
  right=0.66in,
  columnsep=0.25in
]{geometry}
\usepackage[T1]{fontenc}
\usepackage{lmodern}
\usepackage{microtype}
\usepackage{amsmath,amssymb,mathtools}
\usepackage{booktabs,tabularx,array}
\usepackage{enumitem}
\usepackage{balance}
\usepackage{xcolor}
\usepackage{fancyhdr}
\usepackage{titlesec}
\usepackage{hyperref}

\definecolor{Spectrum}{HTML}{3B3F8C}
\definecolor{Ink}{HTML}{171922}
\definecolor{Quiet}{HTML}{626779}
\hypersetup{
  colorlinks=true,
  linkcolor=Spectrum,
  urlcolor=Spectrum,
  pdftitle={Spectrum Beta 1: Probability-Weighted Nametime},
  pdfauthor={HyperBEAM contributors}
}

\setlength{\parindent}{0pt}
\setlength{\parskip}{3pt plus 1pt}
\setlength{\emergencystretch}{2em}
\raggedbottom
\widowpenalty=10000
\clubpenalty=10000
\displaywidowpenalty=10000
\setlist{nosep,leftmargin=1.25em}
\setcounter{secnumdepth}{3}
\setcounter{tocdepth}{2}
\titleformat{\section}{\Large\bfseries\color{Ink}}{\thesection}{0.6em}{}
\titleformat{\subsection}{\large\bfseries\color{Ink}}{\thesubsection}{0.55em}{}
\titleformat{\subsubsection}{\normalsize\bfseries\color{Ink}}{\thesubsubsection}{0.5em}{}
\titlespacing*{\section}{0pt}{11pt}{5pt}
\titlespacing*{\subsection}{0pt}{8pt}{3pt}
\titlespacing*{\subsubsection}{0pt}{6pt}{2pt}

\pagestyle{fancy}
\fancyhf{}
\fancyfoot[C]{\small\color{Quiet}\thepage}
\renewcommand{\headrulewidth}{0pt}
\fancypagestyle{plain}{
  \fancyhf{}
  \fancyfoot[C]{\small\color{Quiet}\thepage}
  \renewcommand{\headrulewidth}{0pt}
}

\newcommand{\device}[1]{\texttt{\textasciitilde#1}}
\newcommand{\key}[1]{\texttt{#1}}
\newcommand{\R}{\mathbb{R}}
\newcommand{\ceil}[1]{\left\lceil #1\right\rceil}
\newcommand{\floor}[1]{\left\lfloor #1\right\rfloor}
\newcommand{\E}{\mathbb{E}}
\newcommand{\Prob}{\operatorname{Pr}}

\begin{document}

\twocolumn[{
  \begin{minipage}{\textwidth}
    \noindent{\color{Spectrum}\rule{0.10in}{0.72in}}\hspace{0.16in}
    \begin{minipage}[b]{0.84\textwidth}
      {\fontsize{24}{28}\selectfont\bfseries\color{Ink}
        Spectrum Beta 1\par}
      \vspace{3pt}
      {\fontsize{14}{17}\selectfont\color{Quiet}
        Probability-Weighted Nametime for a Permaweb Namespace\par}
      \vspace{5pt}
      {\small\color{Quiet}Protocol note, August 2026}
    \end{minipage}

    \vspace{0.18in}
    \hrule
    \vspace{0.10in}
    \textbf{Abstract.}
    Spectrum is a single-market naming protocol. A registration transaction
    pays the Arweave network reward and receives the greatest whole number of
    lease blocks that payment can buy. A fixed Markov model assigns every exact
    name a probability weight; active leases consume that probability mass
    over time. A convex bonding curve prices the integral increase in occupied
    mass, and known expiries divide a quote into constant-occupancy time slabs.
    \device{markov@1.0} defines the statistical model,
    \device{probability-time@1.0} defines the pure pricing function, and
    \device{spectrum@1.0} defines the scheduled registry and lease state
    machine. This note specifies their mathematics, mechanism properties,
    message interfaces, state invariants, and deployment model.
    \vspace{0.16in}
  \end{minipage}
}]

\section{Protocol overview}

\subsection{Design objective}

The allocated resource is \emph{exact-name probability mass over block time}.
Let a fixed distribution assign probability $p(n)$ to exact name $n$. A lease
for $n$ lasting $d$ blocks consumes $p(n)d$ probability-blocks. This makes
scarcity continuous: every valid name uses the same market and differs only in
its weight. There is no protocol branch between ``free'', fixed-price, and
auctioned names.

The beta has four design constraints:

\begin{enumerate}
  \item one deterministic price for every node evaluating the same schedule;
  \item a fixed capacity measure that leases release automatically at expiry;
  \item whole-block grants that never exceed what the payment affords; and
  \item composable devices whose statistical model and pricing function can be
        replaced without changing the registry state machine.
\end{enumerate}

The mechanism prices leases, not ownership in perpetuity. A purchase may fund
an unregistered name or extend a held name. The process does not collect a
separate protocol fee in beta: the authenticated Arweave transaction reward is
the payment.

\subsection{Notation and units}

\begin{tabularx}{\columnwidth}{@{}>{$}l<{$}X@{}}
\toprule
n & exact normalized name; a finite sequence of symbols \tabularnewline
M & fixed Markov model and its versioned parameters \tabularnewline
p_M(n) & model probability of generating exactly $n$, including termination \tabularnewline
w(n) & retained lease weight; by default $p_M(n)$ \tabularnewline
h & Arweave block height supplied by the process schedule \tabularnewline
u(h) & active probability mass immediately at height $h$ \tabularnewline
t & reference occupancy, with $0<t<1$ \tabularnewline
k & marginal winston per probability-block at $u=t$ \tabularnewline
g & grace ratio, represented in process state as integer basis points \tabularnewline
d_i & duration in blocks of a constant-occupancy time slab \tabularnewline
\bottomrule
\end{tabularx}

All payments and returned prices are integer winston. One AR is $10^{12}$
winston. Time is integer Arweave block height. Float calculations are
deterministic within the HyperBEAM execution environment; integer rounding at
the protocol boundary is explicit.

\section{Statistical model}

\subsection{Exact-name probability}

\subsubsection{Tokenization and training}

\device{markov@1.0} is a variable-backoff character Markov model. Its core
vocabulary is the 128 ASCII single-byte symbols plus the terminal token
\key{|end|}. The context is prefixed by $q$ copies of \key{|start|}, where
$q\geq 0$ is the model order. Start and end are model tokens, not bytes emitted
in the generated binary.

Training one binary $x_1\ldots x_L$ appends the event sequence
\[
  x_1,\ldots,x_L,\mathtt{|end|}.
\]
For every event, the model increments its count after the full context of at
most $q$ preceding tokens and after every suffix of that context down to the
empty context. Repeated calls add counts to the existing model. Training is
therefore associative over a fixed order and tokenizer: corpus partitioning
does not change the final counts.

The beta namespace corpus is normalized before training to lowercase ASCII
\key{a-z}, \key{0-9}, and \key{-}. The selected model order is $q=4$. This
normalization is a deployment policy; the Markov device itself accepts general
ASCII binaries.

\subsubsection{Hierarchical interpolation}

Write $C(c,a)$ for the count of output token $a$ after context $c$, and
$C(c)=\sum_a C(c,a)$. At the empty context, the distribution is empirical:
\[
  P(a\mid\epsilon)=\frac{C(\epsilon,a)}{C(\epsilon)}.
\]
For nonempty context $c$, let $s(c)$ remove its first token. The device uses
fixed hierarchical interpolation
\[
  P(a\mid c)=
  \frac{C(c,a)+\lambda P(a\mid s(c))}{C(c)+\lambda},
  \qquad \lambda=129.
\]
The constant is the number of possible output tokens: 128 ASCII symbols and
\key{|end|}. An unseen higher-order context therefore backs off exactly to its
suffix distribution. A symbol absent from the entire training corpus retains
zero probability; this is interpolation over observed support, not add-one
smoothing.

\subsubsection{Likelihood and information measures}\label{sec:information}

For $n=x_1\ldots x_L$, let $c_i$ be the suffix of length at most $q$ preceding
event $x_i$, including start tokens before the first byte. Let
$x_{L+1}=\mathtt{|end|}$. The exact-name likelihood used for pricing is
\[
  p_M(n)=\prod_{i=1}^{L+1} P(x_i\mid c_i).
\]
Including end is essential: it scores not only the spelling but the choice to
stop at that length. The device can omit the terminal factor for exploratory
prefix scoring, but \device{probability-time@1.0} always requests
\key{include-end=true}.

The likelihood endpoint can return a float or the exact product as an integer
pair $(N,D)$ before division. Derived endpoints are
\begin{align*}
  S(n) &= -\log_2 p_M(n) && \text{total surprisal in bits},\\
  \bar S(n) &= \frac{S(n)}{E} && \text{mean surprisal per scored event},\\
  \operatorname{PP}(n) &= 2^{\bar S(n)} && \text{perplexity},
\end{align*}
where $E=L+1$ when end is included. Pricing uses $p_M(n)$ directly, not its
logarithm. Surprisal is the logarithmic effect size and is useful for analysis,
but it does not preserve probability mass under summation.

\subsubsection{The natural occurrence knee}

If $N$ independent names are drawn from $M$, exact name $n$ appears at least
once with probability
\[
  1-(1-p_M(n))^N.
\]
At $p_M(n)=1/N$, this tends to $1-e^{-1}\approx0.6321$. Thus $1/N$ is a
descriptive occurrence knee: above it, exact names are likely to have appeared;
below it, they are increasingly unlikely. It is not a curve parameter in
Spectrum.

The expected probability mass covered by $N$ draws is
\[
  \E[U_N]=\sum_n p_M(n)\left[1-(1-p_M(n))^N\right].
\]
This identity explains why a large count of registered strings need not imply
occupancy near one. The total depends on which probability mass those strings
cover, not on the count alone.

\subsubsection{Deterministic generation}

Generation samples the same interpolated distribution used by likelihood.
For exact integer sampling, let the empty-context outcome weights be $W_0(a)$
with total $T_0$. For each successively longer context $c_j$,
\begin{align*}
  W_j(a) &= C(c_j,a)T_{j-1}+129W_{j-1}(a),\\
  T_j &= \left(C(c_j)+129\right)T_{j-1}.
\end{align*}
These integers represent the interpolated probabilities with a common
denominator. Outcomes are ordered by their bytewise symbol IDs.

Given binary seed $z$ and draw counter $i$, the normative pseudorandom integer
is unsigned $\operatorname{SHA256}(z\mathbin\Vert\operatorname{uint64be}(i))$.
Rejection sampling removes modulo bias before selecting an outcome by cumulative
weight. Generation stops on \key{|end|}, or returns a seed and next counter when
an optional output-length limit is reached. Passing that continuation with the
generated body and a larger limit produces exactly the same output as one call
with the larger limit. Consensus execution must supply a seed; seedless calls
may use runtime entropy.

\subsection{Probability-weighted occupancy}

Each lease records a positive weight. The default is
\[
  w(n)=p_M(n).
\]
A process may configure a weighting device $W$ and instead retain
$w(n)=W(p_M(n),n,M)$. The weighting seam is explicit so that a later protocol
can change the economic measure without changing Markov semantics. Beta 1 uses
the identity function.

For lease $i$, let $d_i$ be its deadline and $w_i$ its retained weight. Active
occupancy at block $h$ is
\[
  u(h)=\sum_{i:h<d_i}w_i.
\]
The grace window does not consume occupancy. A lease in grace retains its
resolver right against takeover, but its priced nametime has ended.

The weight is calculated when a free name is acquired and is retained for all
extensions while that lease remains held. Model replacement or retraining
therefore cannot rewrite an existing lease's resource obligation. A name
reacquired after expiry receives a new weight from the model then in force.
The invariant is
\[
  0\leq u(h)<1.
\]

\subsection{Bonding curve}

\subsubsection{Marginal rate}

Let $t\in(0,1)$ be the reference occupancy and $k>0$ the marginal price at
that point. Define
\[
  A=k\frac{1-t}{t}
  \qquad\text{and}\qquad
  r(u)=A\frac{u}{1-u}.
\]
Then
\[
  r(0)=0,\qquad r(t)=k,\qquad
  \lim_{u\to1^-}r(u)=\infty.
\]
The rate is increasing and convex:
\[
  r'(u)=\frac{A}{(1-u)^2}>0,
  \qquad
  r''(u)=\frac{2A}{(1-u)^3}>0.
\]
The asymptote is a capacity boundary, not a claim that observed namespaces
normally approach total modeled mass.

\subsubsection{Finite-name price per block}

Adding a name of weight $w$ moves occupancy from $u$ to $u+w$. Its price for
one block is the integral of the marginal curve:
\begin{align*}
  c(u,w)
    &=\int_u^{u+w}r(x)\,dx\\
    &=A\left[
      \log\left(\frac{1-u}{1-u-w}\right)-w
    \right],
  \qquad u+w<1.
\end{align*}
This charges the name for every infinitesimal unit of mass it adds, including
its own movement along the curve. Consequently, even at $u=0$, a finite
positive $w$ has positive cost. For small $w$,
\[
  c(u,w)=A\left[
    \frac{u}{1-u}w+
    \frac{w^2}{2(1-u)^2}+O(w^3)
  \right].
\]
The cost is increasing in both $u$ and $w$, convex in the added mass, and
diverges as $u+w\to1^-$. A quote that would reach or cross one is invalid.

For numerical stability, the implementation evaluates the logarithm directly
except when $w/(1-u)<10^{-4}$. In that range it sums the stable tail of
$-\log(1-x)$ after the linear term. This avoids cancellation between two
nearly equal floats without changing the equation being implemented.

\subsection{Pricing over nametime}

\subsubsection{Expiry slabs}

Suppose a quote begins at block $a$ and ends at $b$. Existing deadlines
partition $[a,b)$ into slabs $j$, each with duration $d_j$ and constant
occupancy $u_j$. The unrounded price is
\[
  C(n;a,b)=\sum_j d_j\,c(u_j,w(n)).
\]
Future expiries reduce later slab occupancy automatically. Spectrum therefore
prices the actual occupancy path already committed in state; it does not use a
lagging demand estimator or a separate time-varying auction clock.

For a free name, $a$ is the current scheduled block height. For a held name,
the requested extension begins at its current deadline. Equivalently, the
pricing request supplies the lease's remaining live duration and starts at
current height plus that duration. Grace time is excluded from both the start
and the occupancy slabs.

\subsubsection{Integer boundary}

The price of $B$ blocks is
\[
  \operatorname{price}(B)=\ceil{C(B)}\quad\text{winston}.
\]
Given payment $P$, the grant is
\[
  \operatorname{blocks}(P)=
  \max\{B\in\mathbb Z_{\geq0}:\ceil{C(B)}\leq P\}.
\]
The implementation first divides through the successive constant rates, then
checks the normative rounded price and adjusts by whole blocks. Thus no float
boundary can grant an unaffordable block, and no affordable whole block is
withheld.

\section{Mechanism design}

\subsection{One market}

Every accepted purchase invokes the same rule. A fixed Arweave reward buys
more time for lower-weight names and less time for higher-weight names. If the
minimum network reward buys a useful duration, no additional payment mechanism
is needed; the same transaction both schedules the intent and supplies its
price. There is no threshold at which the protocol switches to a distinct
auction.

The shared state variable is occupied probability mass, not registration
count. A purchase raises occupancy only for its granted interval. Expiry lowers
occupancy without an explicit release message. The curve prices concurrent
scarcity, while the slab integral makes later time cheaper when already-known
leases end.

\subsection{Calibration}

The parameters have separate meanings:
\begin{itemize}
  \item $t$ selects the occupancy at which the operator wants a directly
        interpretable reference rate;
  \item $k$ sets that marginal rate in winston per probability-block;
  \item the corpus, normalization, order, and model fix relative name weights;
  \item $g$ controls recovery time after paid nametime, not price or occupancy.
\end{itemize}

Changing $k$ scales every quote linearly. Changing $t$ changes the same global
scale factor when $k$ is held fixed; the shape as a function of $u$ remains
$u/(1-u)$. These are state parameters and can be finalized at deployment
without changing device code.

\subsection{Lease and grace}

Let a purchase at height $h$ buy $B$ blocks. For a free name,
\[
  \text{deadline}=h+B,
  \qquad
  \text{grace}=h+(1+g)B.
\]
For a held name, the deadline extends from its prior deadline and grace extends
from its prior grace:
\[
  d'=d+B,
  \qquad
  q'=q+(1+g)B.
\]
Carrying grace forward prevents a small extension from deleting recovery time
already purchased. During $[d,q)$ the name resolves to the configured grace
notice. At and after $q$ it is free for a new acquisition. An extension during
grace buys forward from $d$, not from the wall-clock height; lapsed paid time is
not silently restored.

The initial beta issuance grants every imported lease the same term $T_0$ from
height $h_0$:
\[
  d=h_0+T_0,
  \qquad
  q=d+gT_0.
\]
The deployment target of 200 nominal years uses $T_0=52{,}560{,}000$ blocks at
262,800 blocks per nominal year.

\subsection{Resolver right}

A free-name purchase may supply the resolver value. Further payments while the
lease is held extend it but cannot replace that value. Any party may fund an
extension. Beta 1 therefore specifies a held resolver right, not an on-chain
owner identity or transfer mechanism. Ownership and authenticated value updates
are outside this mechanism version.

\section{Device architecture}

\subsection{Composition}

The deployed AO process composes these devices:

\begingroup
\scriptsize
\begin{tabularx}{\columnwidth}{@{}p{0.49\columnwidth}X@{}}
\toprule
Device & Responsibility \\
\midrule
\device{arweave-scheduler@1.0} & Orders base-layer data-free Arweave
transactions in \key{all} mode and attaches block height. \\
\device{spectrum@1.0} & Advances the deterministic clock, validates purchase
shape and authenticated payment, maintains leases, and resolves names. \\
\device{probability-time@1.0} & Computes exact-name weight, occupancy slabs,
bonding-curve price, and its whole-block inverse. \\
\device{markov@1.0} & Stores the fixed statistical model and answers exact
likelihood requests. \\
\bottomrule
\end{tabularx}
\endgroup

The composition is message-based. Spectrum does not call Erlang implementation
functions directly. It resolves a configured \key{pricing-device}; that device
resolves the configured \key{probability-device}. Implementation names may be
device IDs or trusted device names according to normal AO-Core loading rules.

\subsection{Message interfaces}

\device{markov@1.0} exposes these protocol keys:
\begin{itemize}
  \item \key{/train} reads one binary or a list of independent binaries from
        request \key{target}, defaulting to \key{body}, and replaces only the
        \key{model} state with updated counts;
  \item \key{/likelihood} returns float probability by default or exact
        \key{numerator}, \key{denominator}, and \key{events} in integer mode;
  \item \key{/surprisal}, \key{/mean-surprisal}, and \key{/perplexity} derive
        the information measures in Section~\ref{sec:information} from exact
        likelihood; and
  \item \key{/generate} returns \key{body}, \key{seed}, and either a resumable
        \key{continues} state or false on terminal output.
\end{itemize}
The effective order is read from request, base, then existing model, defaulting
to two. Training refuses an order mismatch. Other keys use
\device{message@1.0}.

\device{probability-time@1.0} is a pure quote surface. Request
\key{/price=B} returns the ceil-rounded winston cost for $B$ blocks. Request
\key{/blocks=P} returns the greatest affordable duration and the retained
\key{pricing/weight}. Both require \key{name}; \key{token} defaults to AR; and
\key{duration} locates the interval after already-paid time.

\device{spectrum@1.0} exposes the same quote keys after supplying trusted
remaining duration from registry state. Scheduled \key{/compute} applies an
assignment. A name key resolves through the lease table, optionally loading its
resolver value. Snapshot and normalize preserve the ordinary process state
contract.

\subsection{Process state}

The pricing and registry parameters are read from process state:

\begingroup
\scriptsize
\begin{tabularx}{\columnwidth}{@{}p{0.48\columnwidth}X@{}}
\toprule
Key & Meaning \\
\midrule
\key{pricing-device} & Pricing implementation; beta uses
\device{probability-time@1.0}. \\
\key{probability-device} & Probability implementation; beta uses
\device{markov@1.0}. \\
\key{model} & Trained Markov counts and parameters. \\
\key{target-occupancy} & $t$, strictly between zero and one. \\
\key{price-at-target} & $k$, positive winston per probability-block. \\
\key{weighting-device} & Optional replacement $W$; absent in beta. \\
\key{spectrum-height} & Height of the last applied scheduled assignment. \\
\key{grace-factor} & Nonnegative grace ratio in basis points. \\
\key{grace-notice} & Resolver value returned during grace. \\
\key{names} & Map from exact name to lease record. \\
\bottomrule
\end{tabularx}
\endgroup

Each lease record has the form
\[
\{\key{value},\key{deadline},\key{grace},
  \key{pricing}:\{\key{weight}\}\}.
\]
The pricing metadata is opaque to Spectrum except that the pricing device
receives it back through base state. \device{probability-time@1.0} requires a
positive retained weight.

\subsection{Initialization}

Before the process transaction is mined, deployment prepares two signed
\device{ans104@1.0} messages: the trained \key{model} and the pre-issued
\key{names} map. Each prepared record contains its resolver value, exact-name
weight, deadline, and grace height. Total imported occupancy must remain below
one.

The base-layer process transaction remains data-free. It carries scalar
configuration tags plus \key{model+link} and \key{names+link}, whose values are
the signed message IDs. Decoding those tags yields native AO-Core links, so the
devices read ordinary \key{model} and \key{names} state and resolve the large
children lazily through the store. No \key{initial-namespace} wrapper or custom
initialization callback is required. The model and namespace can be published
as bundled data items without flattening their contents into size-limited
Arweave tags.

After the process is created, the \key{all}-mode scheduler admits only real
base-layer data-free transactions to the purchase sequence.
Section~\ref{sec:deployment} records the deployment boundary.

\subsection{Scheduled execution}

The process clock is the \key{block-height} on an assignment produced by
\device{arweave-scheduler@1.0} in \key{all} mode. Spectrum never reads a live
network tip during compute. Every scheduled data-free transaction advances
\key{spectrum-height}; unrelated transactions otherwise take a constant path
and do not scan or mutate the namespace.

A purchase is selected by the scheduled transaction's \key{path=purchase} tag.
The payment is not read from a same-named tag. Spectrum extracts
\key{field-reward} from the transaction's verified \device{tx@1.0} commitment.
The transaction target remains empty, so the reward is the ordinary Arweave
network reward. A malformed, unauthenticated, unaffordable, or otherwise invalid
purchase is ignored. It cannot fail a public \key{all}-mode process slot.

On success, Spectrum asks \key{/blocks} with payment, name, and trusted
remaining duration. A structured pricing response returns both the integer
grant and \key{pricing/weight}. Spectrum updates the deadline and grace, retains
the old value for a held lease, and stores the new pricing metadata atomically.

\subsection{Reads}

At height $h$, a name reads as follows:
\[
\begin{cases}
\key{value}, & h<\key{deadline},\\
\key{grace-notice}, & \key{deadline}\leq h<\key{grace},\\
\key{not-found}, & h\geq\key{grace}.
\end{cases}
\]
A request may ask to dereference the returned value. Historical reads may name
a height or date. Date resolution bisects immutable Arweave block timestamps;
an unavailable probe fails the requested historical read rather than silently
substituting the process's current height.

\subsection{Failure and determinism boundaries}

The following checks are normative:
\begin{itemize}
  \item model probabilities, transformed weights, $t$, and $k$ are positive;
  \item $t<1$, current occupancy $<1$, and quoted occupancy plus weight $<1$;
  \item amounts, heights, and durations are nonnegative integers;
  \item only AR is accepted in beta;
  \item linked state must resolve to the expected message shape;
  \item price rounds up; purchased duration is the greatest affordable whole
        block; and
  \item invalid public input is a no-op, not a divergent process error.
\end{itemize}

The fixed model, retained weights, schedule-derived height, deterministic float
execution, and explicit integer boundary together make a quote reproducible
from process state. Node-local cache layout and fetch timing do not enter the
economic result.

\balance
\section{Deployment and audit}\label{sec:deployment}

The beta deployment has two phases.

\begin{enumerate}
  \item \textbf{State construction.} Normalize the root namespace, train the
        order-4 model, calculate every \key{include-end=true} weight, and sign
        separate bundled AO messages for the model and pre-issued records.
        Create a data-free process transaction carrying their IDs as
        \key{model+link} and \key{names+link}, with initial height, curve
        parameters, and grace factor in scalar state tags.
  \item \textbf{Live scheduling.} Start the registry under
        \device{arweave-scheduler@1.0} \key{mode=all}. Thereafter, only
        base-layer data-free Arweave transactions enter the schedule. Purchase
        rewards are taken from their transaction commitments.
\end{enumerate}

The root source for the beta is the namespace approximating manifest\\
\key{fQXYPE9MAcfI1wV2CwJ3sJIhgT9btBOlYFOKFDGhAs0}. The deployment artifact must
record the exact source, normalization, model order, initial height, nametime,
$t$, $k$, grace factor, device implementation IDs, and content IDs of all
bundled state messages. These values define the reproducible genesis state.

No beta parameter should be inferred from a node's wall clock, chain tip, or
local environment. No namespace artifact should be flattened into transaction
tags. Before publication, the complete composed process is evaluated through
packaged AO-Core messages: initialization, pricing, acquisition, extension,
grace, expiry, and resolution.

\subsection{Scope and limitations}

Beta 1 intentionally fixes several choices:

\begin{itemize}
  \item The statistical model captures string-form regularity, not trademark,
        semantic, social, or application value.
  \item The model is fixed for the life of a retained lease. A future model
        policy must define version transitions explicitly.
  \item The curve prices occupied probability mass. It does not estimate recent
        demand, clear competing simultaneous bids, or pay proceeds to a seller.
  \item The registry supplies lease and resolver semantics, not an ownership or
        transfer protocol.
  \item Occupancy quote construction currently walks lease state to collect
        active weights and future deadlines. Scaling this read path is an
        implementation concern that must preserve the same slab result.
  \item The model's calibration depends on the declared initial corpus and
        normalization. A biased corpus yields biased weights even when the
        arithmetic is exact.
\end{itemize}

These are explicit protocol boundaries, not hidden secondary mechanisms. The
device seams permit later replacement of the probability model or weight
transform while keeping lease accounting and scheduled payment authentication
separate.

\subsection{Summary}

Spectrum Beta 1 treats exact-name probability over time as the scarce resource.
A fixed Markov model supplies normalized weights; active retained weights form
occupancy; a convex integral curve prices the finite mass added by a lease; and
expiry slabs price the actual future occupancy path. An authenticated Arweave
reward maps to the greatest affordable whole-block duration. The result is one
continuous mechanism for all names, with automatic capacity release and no
separate auction path.

The implementation mirrors that decomposition. Markov owns statistical
meaning, probability-time owns pricing, Spectrum owns the scheduled registry,
and the Arweave scheduler owns ordering and height. Their only shared contract
is structured AO messages and explicit state. That separation is the central
audit property of the beta.

\end{document}
